\section{Smart Home Alarm Implementation}

\subsection{Actor Responsibilities}
The smart home alarm is split into small typed actors with a clear ownership model.
The control unit owns the alarm state machine, the configured PIN, and the delay timers.
The siren actor owns only its on/off state.
The keypad actor collects local key presses and forwards PIN submissions without validation.
The sensor actor is reusable and forwards activations for motion and door/window sensors.

\subsection{Typed Message Protocols}
Each actor exposes a dedicated typed protocol.
The control unit accepts PIN submissions, sensor activations, and internal timeout messages.
The keypad accepts key presses and direct PIN submissions.
The sensor accepts activation messages.
The siren accepts activate, deactivate, and status-query messages for tests and demo code.

\subsection{State Machine and Timers}
The control unit state machine follows five states: \textit{DISARMED}, \textit{EXIT\_DELAY}, \textit{ARMED}, \textit{ENTRY\_DELAY}, and \textit{ALARM}.
The configured PIN arms and disarms the system according to the current state.
Entering \textit{EXIT\_DELAY} starts a single exit-delay timer, and entering \textit{ENTRY\_DELAY} starts a single entry-delay timer.
Timeout expiration is delivered as a normal actor message, and obsolete timers are cancelled when the actor leaves the timed state.

\subsection{Stale Timeouts}
Timeout messages are ignored unless they match the current state.
This keeps the behavior deterministic even if a timeout message arrives after a state change.
The tests exercise this case directly by injecting timeout messages into the control unit protocol.

\subsection{Siren Integration}
The siren is activated only when the control unit enters \textit{ALARM}.
It is deactivated when the control unit returns to \textit{DISARMED}.
Repeated activate and deactivate messages are idempotent.

\subsection{Shared State}
The implementation does not use shared mutable state.
Each actor owns its own state and communicates exclusively through message passing.
This makes the behavior easy to test and keeps the demo deterministic.

\subsection{Build and Run}
The project uses Java 21 and Maven.
Build the project with \texttt{mvn -f assignment-3/shas/pom.xml compile}.
Run the automated tests with \texttt{mvn -f assignment-3/shas/pom.xml test}.
Run the demonstration with \texttt{mvn -f assignment-3/shas/pom.xml exec:java}.

\subsection{Design Decisions}
The implementation favors small actors over a monolithic controller.
Configuration is loaded from the standard Typesafe Config mechanism so tests can override values without editing production files.
The runnable demo uses short timer overrides to show the whole scenario quickly while leaving the production defaults unchanged.
